\documentclass{article}

% NeurIPS 2025 preprint style (as required)
\usepackage[preprint]{neurips_2025}

\usepackage[utf8]{inputenc}
\usepackage[T1]{fontenc}
\usepackage{hyperref}
\usepackage{url}
\usepackage{booktabs}
\usepackage{graphicx}

\title{Agentic Workflows for Innlandet:\\
Innlandet Public Announcements Digest}

\author{
  Group X \\
  University of Inland Norway \\
  \texttt{your.emails@inn.no}
}

\begin{document}

\maketitle

\begin{abstract}
Residents, international students, and volunteers in the Innlandet region regularly miss official public announcements because they are scattered across agency websites, written in formal Norwegian, and not aggregated. This project builds an agentic harness that fetches official RSS feeds (starting with Statsforvalteren i Innlandet), filters citizen-relevant items, and produces a structured plain-language digest whose correctness can be mechanically checked. The walking skeleton demonstrates an end-to-end multi-stage pipeline that uses tools and runs autonomously.
\end{abstract}

\section{Problem Description}
Innlandet’s public sector continuously publishes announcements, hearings, roadwork notices, environmental decisions, and service updates. These appear on different sites (Statsforvalteren, fylkeskommune, municipalities, NRK) and are easy to miss. The recurring need is a single, regularly updated, accessible digest of citizen-facing public information for residents, international students at the University of Inland Norway, and local voluntary organisations. The workflow produces a structured list of relevant announcements containing title, date, link, source, category, and a short plain-language summary (Norwegian, with optional English). No personal data is processed; only public RSS and openly published documents are used.

\section{Definition of Success}
Correct output is a JSON (and accompanying Markdown) document that:
\begin{itemize}
  \item contains at least one item when the source feed has relevant announcements;
  \item every item has non-empty \texttt{title}, \texttt{link} (valid HTTP URL), \texttt{date}, and \texttt{plain\_summary} ($\geq 15$ characters);
  \item the number of items matches the reported count;
  \item relevant items mention local or public-service concepts (hearings, plans, roads, environment, education, housing, etc.).
\end{itemize}
A separate checker script validates these properties automatically. In the full system an evaluation harness will run a fixed test set of feeds and expected outcomes.

\section{Architecture}
The harness is organised in three stages:
\begin{enumerate}
  \item \textbf{Fetch} -- download the official RSS feed using HTTP (tool use).
  \item \textbf{Parse \& filter} -- extract structured items and retain those matching a documented relevance lexicon (context engineering).
  \item \textbf{Summarise \& format} -- produce plain-language summaries and write checkable JSON + Markdown.
\end{enumerate}
Context is supplied via files in \texttt{context/} that state the role, constraints, and relevance criteria. In the walking skeleton the summarisation stage is deterministic; the full system will replace it with an LLM agent stage while keeping the same interfaces. At least one stage (the complete pipeline) runs without a person driving each step via \texttt{run.sh} or \texttt{python scripts/run\_digest.py}.

\section{Schedule}
Working backwards from the demonstration (weeks 46--47) and white-paper deadline (3 December 2026):
\begin{itemize}
  \item 6 November: Compulsory 2 report and walking skeleton submitted.
  \item 10--16 November: Add second source (NRK Innlandet RSS), improve classification, introduce LLM summarisation stage with fixed prompt.
  \item 17--23 November: Build evaluation harness with typical, boundary, and rejection cases; implement basic safe-operation controls (token budget, kill switch).
  \item Weeks 46--47: Live demonstration; address feedback.
  \item 24 November--2 December: Stress testing (two perturbation classes), regression evidence, final white paper.
  \item 3 December 11:59: White paper and tagged repository submitted.
\end{itemize}

\section{Risk Assessment}
\begin{itemize}
  \item \textbf{Source instability}: RSS structure or availability changes. Mitigation: pin to official feeds, graceful degradation, fallback to cached sample, monitor feed health in evaluation harness.
  \item \textbf{LLM hallucination or format drift}: Agent invents links or alters facts. Mitigation: deterministic extraction of title/link/date; summarisation constrained by context files; automated schema checks; human-in-the-loop gate before any external publication.
  \item \textbf{Over- or under-filtering}: Keyword filter misses relevant items or includes noise. Mitigation: expand lexicon from real items, later replace with agent classification evaluated on a labelled test set.
  \item \textbf{Cost or rate limits}: Full system exceeds budget. Mitigation: explicit token budget per run and hard kill switch.
\end{itemize}
Repository: \url{https://github.com/YOUR_GROUP/innlandet-digest} (to be created/updated with the skeleton).

\section*{Appendix}
\subsection*{Prompt logs (excerpts)}
\emph{(Maximum two pages -- to be filled with representative agent interactions during development.)}

Example skeleton-stage prompt used for future LLM summarisation:
\begin{quote}
You are a local information assistant for Innlandet. Summarise the following official announcement in plain Norwegian (max 80 words). Do not invent facts. Keep the original date and link.\\
Title: \ldots\\
Description: \ldots
\end{quote}

\subsection*{Meeting reports}
\emph{(Evidence of iterative reasoning. Weekly meetings strongly encouraged.)}

\begin{itemize}
  \item 5 November 2026: Agreed on problem domain (public announcements accessibility). Selected Statsforvalteren RSS as first source. Defined success criteria and technical minimum.
  \item 6 November 2026: Implemented walking skeleton, verified end-to-end run and automated checker. Drafted Compulsory 2 report.
\end{itemize}

\end{document}
